\chapter{The zeno Command}\label{app-commands}

Everything Zeno does is done by one command, \verb|zeno|, followed by
what to do and usually a file. This appendix lists every command and
every option of version 0.1. \verb|zeno --help| prints the same list.

\section{Commands}

\begin{longtable}{lP{8.6cm}}
\toprule
Command & What it does \\
\midrule
\endhead
\verb|zeno run f.zeno| & check the program, then run it. Words after
  the file are the program's own arguments
  (\hyperref[ch-reading-writing]{Reading and Writing}) \\
\verb|zeno check f.zeno| & check the program for mistakes, and run
  nothing \\
\verb|zeno build f.zeno| & compile to C and then to a program that runs
  by itself (\hyperref[ch-compiling]{Compiling}). The C is kept beside
  it \\
\verb|zeno test| & run the \verb|test "name" { }| blocks of a file, of
  every file in a folder, or of this folder when nothing is named
  (\hyperref[ch-testing]{Testing}). \verb|--only words| runs only the
  tests with those words in their names \\
\verb|zeno format f.zeno| & lay files out the one way Zeno lays them
  out; several files, or a folder, can be named
  (\hyperref[ch-tools]{The Tools}) \\
\verb|zeno convert f.py| & write a file of Python again as Zeno, as
  \verb|f.zeno|, keeping what cannot be written in an
  \verb|inline python| block, and say what was done with each piece
  (\hyperref[ch-convert]{A File of Python Written as Zeno}).
  \verb|--to| names the file to write, \verb|--replace| writes over
  one \\
\verb|zeno convert f.c| & the same for a file of C, keeping what
  cannot be written in an \verb|inline c| block
  (\hyperref[ch-convert-c]{A File of C Written as Zeno}) \\
\verb|zeno new name| & make a folder with a new program in it: a
  \verb|main.zeno|, a module with a test, and a \verb|zeno.package|
  (\hyperref[ch-packages]{Packages}) \\
\verb|zeno new --package name| & make a folder with a new package: modules
  for other programs to need, and no \verb|main.zeno| \\
\verb|zeno fetch| & get the packages that the \verb|needs| lines of
  \verb|zeno.package| list, for this folder or a named one \\
\verb|zeno search words| & list the packages of the registry that have
  these words in their name or in what they are for; with no word, all
  of them \\
\verb|zeno publish| & check that the package in this folder is ready
  for the registry, and say how to ask the registry to take it. A
  version after it, \verb|zeno publish v1.2.0|, says which tag \\
\verb|zeno registry accept| & for whoever keeps a registry: take a
  version of a package into a clone of it
  (\hyperref[ch-packages]{Packages}) \\
\verb|zeno tokens f.zeno| & show the tokens the lexer finds
  (\hyperref[ch-inside]{Inside Zeno}) \\
\verb|zeno parse f.zeno| & show the syntax tree the parser builds, with
  brackets around every expression \\
\verb|zeno join f.zeno| & show the one program that the file and every
  file it imports are joined into \\
\verb|zeno bytecode f.zeno| & show the bytecode that \verb|zeno run --vm|
  runs \\
\verb|zeno playground| & write \verb|playground.html|, a page in which
  Zeno programs are written and run in a browser, with nothing installed
  (\hyperref[ch-web]{A Program in a Web Page}) \\
\verb|zeno lsp| & the language server. An editor starts it, to show
  mistakes as you type; you do not run it yourself \\
\verb|zeno --version| & print the version \\
\verb|zeno --help| & print the list of commands and options \\
\bottomrule
\end{longtable}

\section{Options}

An option goes after the command and before the file.

\begin{longtable}{llP{6.6cm}}
\toprule
Option & With & What it does \\
\midrule
\endhead
\verb|--strict| & \verb|run|, \verb|check| & every type must be known
  before the program runs, as a build needs
  (\hyperref[ch-types]{Types, Written and Worked Out}). \verb|build| always works so \\
\verb|--vm| & \verb|run| & compile to bytecode and run that, in place of
  walking the tree. Same output, about twice as fast \\
\verb|--no-collector| & \verb|check| & check as
  \verb|build --no-collector| would, without building \\
\verb|--firmware| & \verb|build| & errors travel as return values, with
  no \verb|setjmp| (\hyperref[ch-firmware]{Firmware}) \\
\verb|--no-collector| & \verb|build| & no garbage collector and no
  \verb|malloc|: every list and object has one owner, and memory comes
  from a fixed pool (\hyperref[ch-no-collector]{Building with No
  Collector}). Turns \verb|--firmware| on \\
\verb|--owned| & \verb|build|, \verb|check|, \verb|run| & one owner for
  every list and object and no collector, as \verb|memory owned| in
  \verb|zeno.package| says; memory grows as the program needs it
  (\hyperref[ch-memory]{Memory}) \\
\verb|--collector| & \verb|build|, \verb|check|, \verb|run| & the garbage
  collector, whatever \verb|zeno.package| says \\
\verb|--release| & \verb|build| & whole numbers wrap round on overflow
  instead of stopping the program; every other check stays \\
\verb|--fast| & \verb|build| & no index check, no check of a float
  divided by zero or of a square root below zero, overflow with no
  defined result (as in C), and \verb|-O3| \\
\verb|--minimal| & \verb|build| & the least C that runs the program:
  no check of any kind, only the parts of the runtime the program
  uses, \verb|print| as \verb|printf|, and \verb|-Os|
  (\hyperref[ch-compiling]{Compiling}; new in 0.5.2) \\
\verb|--board <folder>| & \verb|build| & build for a board in one
  step: the \verb|zeno.board| in the folder names the board's C
  compiler, flags, start-up files and linker script, and the program
  comes out as \verb|name.elf|. Built with no collector. Goes with
  \verb|--minimal|, \verb|--release|, \verb|--fast| and
  \verb|--memory| (\hyperref[ch-compiling]{Compiling}; new in
  0.7) \\
\verb|--llvm| & \verb|build| & write LLVM IR in place of C and hand it
  to \verb|clang|. Goes with the three options around it \\
\verb|--web| & \verb|build| & write a web page that runs the program:
  \verb|name.html| beside it, or a plain page, with the program inside,
  as \verb|name.web.html| (\hyperref[ch-web]{A Program in a Web
  Page}). Goes with no other build option \\
\verb|--target| & \verb|build| & write the program in another language: \verb|js|, \verb|ts|, \verb|html|, \verb|py|, \verb|php|, \verb|go|, \verb|java|, \verb|cs|, \verb|cpp|, \verb|rust| or \verb|wasm| (\hyperref[ch-compiling]{Compiling}; new in 0.5), and since 0.5.2 \verb|lua|, \verb|ruby|, \verb|dart| and \verb|kotlin| \\
\verb|--memory| & \verb|build| & the program reports, as it ends, how
  much memory it used (\hyperref[ch-memory]{Memory}) \\
\verb|--check| & \verb|format| & change nothing; name the files that are
  not laid out, and fail if there are any \\
\verb|--package| & \verb|new| & make a package, not a program \\
\verb|--update| & \verb|fetch| & look every version up again, and not at
  the commits that \verb|zeno.lock| pinned \\
\bottomrule
\end{longtable}

\section{What a command answers}

Every command ends with an \emph{exit status}, the number a shell
script or a build tool reads to know how it went: 0 when all went
well, and 1 when the program had a mistake, an uncaught error, a
failing test, or a file that \verb|format --check| would change.
Messages go to the error stream, and what the program prints goes to
the output stream, so \verb|zeno run f.zeno > out.txt| keeps only what
the program printed.

\section{Telling the C compiler}

A build writes C, and these names can be set when that C is compiled,
with \verb|-D| on the line of the C compiler:

\begin{longtable}{lP{8.6cm}}
\toprule
Name & What it sets \\
\midrule
\endhead
\verb|ZENO_BOARD| & a header that puts a real board under the program,
  in place of the simulated one: \verb|-DZENO_BOARD='"my_board.h"'|
  (\hyperref[ch-firmware]{Firmware}) \\
\verb|ZENO_BOARD_PRINTS| & said by such a header, with
  \verb|#define|, when it brings the four functions that print, warn,
  read and stop in place of the C library
  (\hyperref[ch-firmware]{Firmware}) \\
\verb|ZENO_POOL| & the size of the pool in bytes for a build with no
  collector, 32\,768 unless set: \verb|-DZENO_POOL=25384| \\
\verb|ZENO_TEXT_ROOM| & the size of the text room in bytes, for the same
  build \\
\verb|ZENO_PLAIN_C| & asks GCC or Clang for the checks written in
  plain C, which every other C compiler gets unasked
  (\hyperref[ch-firmware]{Firmware}; new in 0.7) \\
\bottomrule
\end{longtable}

\section{Names in the environment}

Zeno finds the other programs it needs by their usual names. A name
in the environment says where another one is:

\begin{longtable}{lP{8.6cm}}
\toprule
Name & What it names \\
\midrule
\endhead
\verb|CC| & the C compiler of \verb|zeno build| and of
  \verb|import c|, in place of \verb|cc|, \verb|gcc| or \verb|clang| \\
\verb|ZENO_PYTHON|, \verb|ZENO_NODE| & the Python or the Node that
  \verb|import python| or \verb|import node| starts
  (\hyperref[ch-other-languages]{Libraries of Other Languages}) \\
\verb|ZENO_JAVA|, \verb|ZENO_RUBY| & the same for \verb|import java|
  and \verb|import ruby| \\
\verb|ZENO_PHP|, \verb|ZENO_LUA| & and for \verb|import php| and
  \verb|import lua| \\
\verb|ZENO_DOTNET| & the \verb|dotnet| of \verb|import csharp| \\
\bottomrule
\end{longtable}
